\documentclass[12pt, a4paper]{article}

\usepackage[utf8]{inputenc}
\usepackage[spanish]{babel}
\usepackage{geometry}
\usepackage{titlesec}
\usepackage{hyperref}
\usepackage{enumitem}
\usepackage{xcolor}

\geometry{margin=2.5cm}

\definecolor{cabecera}{RGB}{30, 90, 160}

\titleformat{\section}{\large\bfseries\color{cabecera}}{}{0em}{}[\titlerule]

% --------------------------------------------------
\begin{document}

\begin{center}
    {\LARGE \textbf{Reparto de Tareas}}\\[0.4em]
    {\large \textit{Space Invaders — Proyecto de Ingeniería del Software}}\\[0.3em]
    {\small Fecha: \today}
\end{center}

\vspace{1em}
\hrule
\vspace{1.5em}

% --------------------------------------------------
\section{Miembros del equipo Cachopín}

\begin{itemize}[leftmargin=2em]
    \item \textbf{Miembro 1:} Unai Urrutia
    \item \textbf{Miembro 2:} Ekain Ortega
    \item \textbf{Miembro 3:} Eder Lorenzo
    \item \textbf{Miembro 4:} Asier López
\end{itemize}

\vspace{1.5em}

% --------------------------------------------------
\section{Planificación inicial}

Tareas previstas antes de comenzar el desarrollo y su tiempo estimado:

\begin{enumerate}[leftmargin=2em]
    \item Diseño del juego Espacio (MVC) — 8 horas
    \item Pantallas StartFrame y MainFrame — 6 horas
    \item Vista y controlador y métodos update— 12 horas
    \item Implementación de la lógica de juego, disparo y colisiones — 8 horas
    \item Solución de errores en patrón MVC+Observer y el juego — 5 horas
    \item Documentación, diagramas ... — 2 horas
\end{enumerate}

\textbf{Tiempo total estimado:} 41 horas.

\vspace{1.5em}

% --------------------------------------------------
\section{Resumen de tareas realizadas}

Tareas efectivamente completadas y su tiempo real:

\begin{enumerate}[leftmargin=2em]
    \item Entendimiento del patrón MVC e ideamiento de las clases principales del juego — 5 horas
    \item Creación del repositorio y cambios para evitar problemas con archivos de configuración — 1 hora
    \item Desarrollo de las clases principales del modelo — 5 horas
    \item Desarrollo de la lógica de juego, disparo y colisiones — 4 horas
    \item Desarrollo de la vista, el controlador y los métodos update — 3 horas
    \item Solución de los métodos notify y update para inicializar MainFrame correctamente — 4 horas
    \item Solución de errores: colisiones, límites, movimiento — 3 horas
    \item Refinamiento del modelo — 1 hora
    \item Solución de estructura de MainFrame para implementarla con labels — 5 horas
    \item Reestructuración del update y notify para MainFrame, evitando que MainFrame tenga lógica de negocio — 5 horas
    \item Revisión y cambios para implementar correctamente el patrón MVC+Observer — 5 horas

    \item Documentación — 6 horas
    \item Diagramas UML y del patrón MVC+Observer — 3 horas
\end{enumerate}

\textbf{Tiempo total real:} 50 horas.

\vspace{1.5em}

% --------------------------------------------------
\section{Detalle de tareas}

% ---- Tarea 1 ----
\subsection*{Tarea 1 — Entendimiento del patrón MVC e ideamiento de las clases principales del juego}
\begin{description}[leftmargin=2em, style=nextline]
    \item[\textbf{Descripción:}] Estudio del patrón MVC y Observer, diseño conceptual de las clases necesarias para el juego: \texttt{Espacio}, \texttt{Jugador}, \texttt{Enemigo}, \texttt{Naves}, \texttt{Disparo}, y planificación de la arquitectura general. Además, pensamos y miramos en otras prácticas cómo se implementó el patrón MVC+Observer.
    \item[\textbf{Tiempo:}] 5 horas.
    \item[\textbf{Responsables y dedicación:}]~
        \begin{itemize}[leftmargin=1em, nosep]
            \item Unai Urrutia — 2.5 horas
            \item Asier López — 1.5 horas
            \item Eder — 0.5 horas
            \item Ekain Ortega — 0.5 horas
        \end{itemize}
\end{description}

\vspace{0.8em}

% ---- Tarea 2 ----
\subsection*{Tarea 2 — Creación del repositorio y configuración inicial}
\begin{description}[leftmargin=2em, style=nextline]
    \item[\textbf{Descripción:}] Creación del repositorio Git, configuración del \texttt{.gitignore} y ajustes para evitar problemas con archivos de configuración del IDE.
    \item[\textbf{Tiempo:}] 1 hora.
    \item[\textbf{Responsables y dedicación:}]~
        \begin{itemize}[leftmargin=1em, nosep]
            \item Unai Urrutia — 0.3 horas
            \item Asier López — 0.7 horas
            \item Eder — 0 horas
            \item Ekain Ortega — 0 horas
        \end{itemize}
\end{description}

\vspace{0.8em}

% ---- Tarea 3 ----
\subsection*{Tarea 3 — Desarrollo de las clases principales del modelo}
\begin{description}[leftmargin=2em, style=nextline]
    \item[\textbf{Descripción:}] Implementación de las clases \texttt{Naves}, \texttt{Jugador}, \texttt{Enemigo}, \texttt{Disparo} y \texttt{Espacio}. Definición de atributos, constructores y métodos básicos. Se dejaron bastantes métodos de Espacio sin implementar.
    \item[\textbf{Tiempo:}] 5 horas.
    \item[\textbf{Responsables y dedicación:}]~
        \begin{itemize}[leftmargin=1em, nosep]
            \item Unai Urrutia — 2.5 hora
            \item Asier López — 1.5 hora
            \item Eder — 0.5 horas
            \item Ekain Ortega — 0.5 horas
        \end{itemize}
\end{description}

\vspace{0.8em}

% ---- Tarea 4 ----
\subsection*{Tarea 4 — Desarrollo de la lógica de juego, disparo y colisiones}
\begin{description}[leftmargin=2em, style=nextline]
    \item[\textbf{Descripción:}] Implementación de la lógica de disparo, detección de colisiones entre disparos y enemigos, condiciones de victoria y derrota. Todo ello en la clase Espacio.
    \item[\textbf{Tiempo:}] 4 horas.
    \item[\textbf{Responsables y dedicación:}]~
        \begin{itemize}[leftmargin=1em, nosep]
            \item Unai Urrutia — 1.5 horas
            \item Asier López — 1.5 horas
            \item Eder — 0.5 horas
            \item Ekain Ortega — 0.5 horas
        \end{itemize}
\end{description}

\vspace{0.8em}

% ---- Tarea 5 ----
\subsection*{Tarea 5 — Desarrollo de la vista, el controlador y los métodos update}
\begin{description}[leftmargin=2em, style=nextline]
    \item[\textbf{Descripción:}] Desarrollo de \texttt{StartFrame} y \texttt{MainFrame}. Implementación de los controladores de teclado (\texttt{KeyListener}), registro como observadores y métodos \texttt{update} para reaccionar a los cambios del modelo.
    \item[\textbf{Tiempo:}] 3 horas.
    \item[\textbf{Responsables y dedicación:}]~
        \begin{itemize}[leftmargin=1em, nosep]
            \item Unai Urrutia — 1 hora
            \item Asier López — 1 hora
            \item Eder — 0.5 horas
            \item Ekain Ortega — 0.5 horas
        \end{itemize}
\end{description}

\vspace{0.8em}

% ---- Tarea 6 ----
\subsection*{Tarea 6 — Solución de los métodos notify y update para inicializar MainFrame correctamente}
\begin{description}[leftmargin=2em, style=nextline]
    \item[\textbf{Descripción:}] Corrección de un problema que detectamos junto al profesor. Startframe no notificaba a Espacio la inicialización del juego, por lo que no se respetaba el patrón MVC+Observer.
    \item[\textbf{Tiempo:}] 4 horas.
    \item[\textbf{Responsables y dedicación:}]~
        \begin{itemize}[leftmargin=1em, nosep]
            \item Unai Urrutia — 1.5 horas
            \item Asier López — 1.5 horas
            \item Eder — 0.25 horas
            \item Ekain Ortega — 0.25 horas
        \end{itemize}
\end{description}

\vspace{0.8em}

% ---- Tarea 7 ----
\subsection*{Tarea 7 — Solución de errores: colisiones, límites, movimiento}
\begin{description}[leftmargin=2em, style=nextline]
    \item[\textbf{Descripción:}] Tras poder empezar a probar el juego, se detectaron errores relacionados con la detección de colisiones, los límites del tablero y el movimiento del jugador y enemigos. Se solucionaron todos estos errores.
    \item[\textbf{Tiempo:}] 3 horas.
    \item[\textbf{Responsables y dedicación:}]~
        \begin{itemize}[leftmargin=1em, nosep]
            \item Unai Urrutia — 1 hora
            \item Asier López — 1 hora
            \item Eder — 0.5 horas
            \item Ekain Ortega — 0.5 horas
        \end{itemize}
\end{description}

\vspace{0.8em}

% ---- Tarea 8 ----
\subsection*{Tarea 8 — Refinamiento del modelo}
\begin{description}[leftmargin=2em, style=nextline]
    \item[\textbf{Descripción:}] Mejoras en la estructura del modelo: limpieza de código, ajuste de responsabilidades entre clases y optimización de métodos.
    \item[\textbf{Tiempo:}] 1 hora.
    \item[\textbf{Responsables y dedicación:}]~
        \begin{itemize}[leftmargin=1em, nosep]
            \item Unai Urrutia — 0.25 horas
            \item Asier López — 0.25 horas
            \item Eder — 0.25 horas
            \item Ekain Ortega — 0.25 horas

        \end{itemize}
\end{description}

\vspace{0.8em}

% ---- Tarea 9 ----
\subsection*{Tarea 9 — Solución de estructura de MainFrame para implementarla con labels}
\begin{description}[leftmargin=2em, style=nextline]
    \item[\textbf{Descripción:}] Tras el feedback del profesor vimos que MainFrame, y por lo tanto, la parte más importante del juego, no estaba implementada correctamente. Tuvimos que reestructurarlo y borrar mucho código. \texttt{MainFrame} pasó a ser una cuadrícula de \texttt{JLabel}. Además tuvimos que hacer cambios en update y notify, como se detalla en la siguiente tarea.
    \item[\textbf{Tiempo:}] 5 horas.
    \item[\textbf{Responsables y dedicación:}]~
        \begin{itemize}[leftmargin=1em, nosep]
            \item Unai Urrutia — 2 horas
            \item Asier López — 1.5 horas
            \item Eder — 0.5 horas
            \item Ekain Ortega — 1 horas
        \end{itemize}
\end{description}

\vspace{0.8em}

% ---- Tarea 10 ----
\subsection*{Tarea 10 — Reestructuración del update y notify para MainFrame}
\begin{description}[leftmargin=2em, style=nextline]
    \item[\textbf{Descripción:}] Mainframe pasó de recibir un update vacío, a recibir un array diciendo qué arrays pintar y de qué color, y cuáles pintar del color del fondo. También se eliminaron variables de estado en \texttt{MainFrame} (posiciones del jugador, disparo, etc.). Toda la información necesaria para pintar y borrar se pasa a través de las notificaciones del modelo, evitando que la vista contenga lógica de negocio.
    \item[\textbf{Tiempo:}] 5 horas.
    \item[\textbf{Responsables y dedicación:}]~
        \begin{itemize}[leftmargin=1em, nosep]
            \item Unai Urrutia — 2 horas
            \item Asier López — 2 horas
            \item Eder — 0.5 horas
            \item Ekain Ortega — 0.5 horas
        \end{itemize}
\end{description}

\vspace{0.8em}

% ---- Tarea 11 ----
\subsection*{Tarea 11 — Revisión y cambios para implementar correctamente el patrón MVC+Observer}
\begin{description}[leftmargin=2em, style=nextline]
    \item[\textbf{Descripción:}] Revisión general del código para asegurar que se respeta el patrón MVC: la lógica de fin de partida se controla desde el modelo, el controlador no contiene lógica de negocio, y la vista solo reacciona a notificaciones. En concreto, se implementaron de forma diferente ciertas llamadas directas de MainFrame a Espacio, que detectaban si el juego había terminado. 
    Además, en el último momento hemos visto graves fallos que han requerido de cambios importantes en la inicialización de MainFrame.
    \item[\textbf{Tiempo:}] 2.5 horas.
    \item[\textbf{Responsables y dedicación:}]~
        \begin{itemize}[leftmargin=1em, nosep]
            \item Unai Urrutia — 2 horas
            \item Asier López — 2 horas
            \item Eder — 0.50 horas
            \item Ekain Ortega — 0.5 horas
        \end{itemize}
\end{description}

\vspace{0.8em}

% ---- Tarea 12 ----
\subsection*{Tarea 12 — Documentación}
\begin{description}[leftmargin=2em, style=nextline]
    \item[\textbf{Descripción:}] Redacción de la memoria del proyecto, justificación del patrón MVC+Observer, y elaboración de este documento de reparto de tareas.
    \item[\textbf{Tiempo:}] 6 horas.
    \item[\textbf{Responsables y dedicación:}]~
        \begin{itemize}[leftmargin=1em, nosep]
            \item Eder — 1 hora
            \item Ekain Ortega — 1 horas
            \item Unai Urrutia — 2 horas
            \item Asier López — 2 horas
        \end{itemize}
\end{description}

\vspace{0.8em}

% ---- Tarea 13 ----
\subsection*{Tarea 13 — Diagramas UML y del patrón MVC+Observer}
\begin{description}[leftmargin=2em, style=nextline]
    \item[\textbf{Descripción:}] Elaboración de los diagramas UML de clases y de la arquitectura MVC+Observer del proyecto.
    \item[\textbf{Tiempo:}] 3 horas.
    \item[\textbf{Responsables y dedicación:}]~
        \begin{itemize}[leftmargin=1em, nosep]
            \item Unai Urrutia — 1 hora
            \item Eder — 1 horas
            \item Ekain Ortega — 1 horas
            \item Asier López — 0 horas
        \end{itemize}
\end{description}

\end{document}
